% Bagian: counterfactuals
% Bab: introduction
% Subbagian: strict-conditional

\documentclass[../../../include/open-logic-section]{subfiles}

\begin{document}

\olfileid[id]{cnt}{int}{str}

\olsection{Kondisional Ketat}

Lewis memperkenalkan \emph{kondisional ketat}~$\strictif$ dan berpendapat
bahwa kondisional inilah, bukan kondisional material, yang bersesuaian dengan
implikasi. Dalam logika modal aletik, $!A \strictif !B$ dapat didefinisikan
sebagai $\Box(!A \lif !B)$. Jadi, suatu kondisional ketat benar (pada suatu
dunia) jika dan hanya jika kondisional material yang bersesuaian bersifat
niscaya.

Bagaimana kondisional ketat menghadapi paradoks-paradoks kondisional
material? Kondisional ketat dengan anteseden yang salah, maupun kondisional
ketat dengan konsekuen yang benar, mungkin benar atau mungkin salah. Selain
itu, $(!A \strictif !B) \lor (!B \strictif !A)$ tidak valid. Kondisional
ketat $!A \strictif !B$ juga tidak ekuivalen dengan $\lnot !A \lor !B$;
jadi, kondisional itu tidak bersifat fungsional-kebenaran.

Kita mempunyai:
\begin{align}
  !A \strictif !B & \Entails \lnot !A \lor !B\\
  \intertext{tetapi:}
  \lnot !A \lor !B & \Entails/ !A \strictif !B\\
  !B & \Entails/ !A \strictif !B\\
  \lnot !A & \Entails/ !A \strictif !B\\
  \lnot(!A \strictif !B) & \Entails/ !A \land \lnot !B\\
  \intertext{tetapi:}
  !A \land \lnot !B & \Entails \lnot(!A \strictif !B)
  \intertext{Namun, kondisional ketat tetap mendukung modus ponens:}
  !A, !A \strictif !B & \Entails !B
  \intertext{Kondisional ketat memungkinkan aglomerasi:}
  !A \strictif !B,  !A \strictif !C & \Entails !A \strictif (!B \land !C)
  \intertext{Penguatan anteseden berlaku bagi kondisional ketat:}
  !A \strictif !B & \Entails (!A \land !C) \strictif !B
  \intertext{Kondisional ketat juga bersifat transitif:}
  !A \strictif !B, !B \strictif !C & \Entails !A \strictif !C
  \intertext{Akhirnya, kondisional ketat ekuivalen dengan kontraposisinya:}
  !A \strictif !B & \Entails \lnot !B \strictif \lnot !A\\
  \lnot !B \strictif \lnot !A & \Entails !A \strictif !B
\end{align}

\begin{prob}
  Berikan contoh tandingan dalam~\Log{S5} bagi relasi perikutan yang tidak
  berlaku untuk kondisional ketat, yakni bagi:
  \begin{enumerate}
  \item $\lnot p \Entails/ \Box(p \lif q)$
  \item $q \Entails/ \Box(p \lif q)$
  \item $\lnot\Box(p \lif q) \Entails/ p \land \lnot q$
  \item $\Entails/ \Box(p \lif q) \lor \Box(q \lif p)$
  \end{enumerate}
\end{prob}

\begin{prob}
  Tunjukkan bahwa relasi-relasi perikutan yang valid berlaku bagi
  kondisional ketat dengan memberikan bukti-bukti dalam~\Log{S5} bagi:
  \begin{enumerate}
  \item  $\Box(!A \lif !B) \Entails \lnot !A \lor !B$
  \item $!A \land \lnot !B \Entails \lnot\Box(!A \lif !B)$
  \item $!A, \Box(!A \lif !B) \Entails !B$
  \item $\Box(!A \lif !B), \Box(!A \lif !C) \Entails \Box(!A \lif (!B
    \land !C))$
  \item $\Box(!A \lif !B) \Entails \Box((!A \land !C) \lif !B)$
  \item $\Box(!A \lif !B), \Box(!B \lif !C) \Entails \Box(!A
    \lif !C)$
  \item $\Box(!A \lif !B) \Entails \Box(\lnot !B \lif \lnot !A)$
  \item $\Box(\lnot !B \lif \lnot !A) \Entails \Box(!A \lif !B)$
  \end{enumerate}
\end{prob}

Namun, kondisional ketat tetap mempunyai ``paradoksnya'' sendiri. Sebagaimana
kondisional material dengan anteseden salah atau konsekuen benar adalah
benar, kondisional ketat dengan anteseden yang \emph{niscaya} salah atau
konsekuen yang niscaya benar adalah benar. Selain itu, setiap kondisional
ketat yang benar bersifat niscaya benar, dan setiap kondisional ketat yang
salah bersifat niscaya salah. Dengan kata lain, kita mempunyai
\begin{align}
  \Box \lnot !A & \Entails !A \strictif !B\\
  \Box !B & \Entails !A \strictif !B\\
  !A \strictif !B& \Entails \Box(!A \strictif !B)\\
  \lnot(!A \strictif !B) & \Entails \Box\lnot(!A \strictif !B)
\end{align}
Hal-hal ini bukan masalah jika Anda memandang $\strictif$ sebagai
``mengimplikasikan.'' Bagaimanapun, relasi perikutan logis merupakan fakta
matematis sehingga tidak mungkin bersifat kontingen. Namun, hal-hal tersebut
memang menimbulkan persoalan jika Anda hendak menggunakan $\strictif$ sebagai
konektif logis yang dimaksudkan untuk menangkap ``jika \dots, maka \dots,''
terutama dua yang terakhir. Tentulah ada pernyataan ``jika \dots, maka
\dots'' yang benar secara kontingen atau salah secara kontingen---bahkan,
pernyataan semacam itu pada umumnya tidak niscaya maupun mustahil.

\begin{prob}
  Berikan bukti dalam~\Log{S5} bagi:
  \begin{enumerate}
    \item  $\Box \lnot !A \Entails !A \strictif !B$
    \item $!A \strictif !B \Entails \Box(!A \strictif !B)$
    \item $\lnot(!A \strictif !B)  \Entails
      \Box\lnot(!A \strictif !B)$
  \end{enumerate}
  Gunakan definisi $\strictif$ untuk melakukannya.
\end{prob}

\end{document}
